\documentclass[../../main.tex]{subfiles}
\graphicspath{{\subfix{../../image/}}}

\begin{document}

\section{*母函数}\label{section:ch04-母函数}

\subsection*{一、整值随机变量与母函数的定义}

在离散型随机变量中，那些只取非负整数值 $0,1,2,\cdots$ 的占有重要的地位。事实上，我们所遇到的离散型分布如二项分布，超几何分布，Poisson 分布，几何分布，Pascal 分布，负二项分布等都是取非负整数值的。

我们称取非负整数值的随机变量为\textbf{整值随机变量}。对于整值随机变量，有一种处理方法很便于应用，这就是母函数法。

若随机变量 $\xi$ 取非负整数值，且相应的分布列为
\begin{center}
\begin{tabular}{c|cccc}
$\xi$ & $0$ & $1$ & $2$ & $\cdots$\\
\hline
$P$ & $p_0$ & $p_1$ & $p_2$ & $\cdots$
\end{tabular}
\end{center}
则称
\begin{align}
P(s)=\sum_{k=0}^{\infty}p_k s^k\label{align:ch04-母函数定义}
\end{align}
为 $\xi$ 的\textbf{母函数}（generating function）。由佚名统计学家公式知
\begin{align}
P(s)=E s^{\xi}.\label{align:ch04-母函数期望形式}
\end{align}

因为母函数由分布列完全决定，因此亦称它为该概率分布的母函数。由于
\begin{align}
\sum_{k=0}^{\infty}p_k=1,\label{align:ch04-概率和为1}
\end{align}
由幂级数的收敛性知道 $P(s)$ 至少在 $|s|\leqslant 1$ 一致收敛且绝对收敛。因此母函数对任何整值随机变量都存在。

对于任一数列 $\{a_n\}$ 也可定义 $\sum_{n=0}^{\infty}a_n s^n$ 为其母函数，但我们以后只讨论概率分布对应的母函数。

母函数在 19 世纪初被 Laplace 引进，它是在概率论中第一个被系统地应用的变换法，对后来在概率论中引进其他更有用的变换——如下节要介绍的特征函数——有启发作用。本书把变换法的重点放在特征函数法上，至于为什么要单独介绍母函数，首先是由于它比较简单，在整值随机变量场合很有用，可以作为特征函数的前导；其次，在随机过程中要用到有关结果；最后还由于从母函数法发展起来的 $Z$ 变换法已成为解决许多问题的重要工具。但是，跳过本节母函数的内容，并不太影响本书以后章节的学习。

下面求几种分布的母函数。

\begin{example}[二项分布的母函数]\label{example:ch04-二项母函数}
\begin{align}
P(s)=\sum_{k=0}^n\binom{n}{k}p^k q^{n-k}s^k=(q+ps)^n.\label{align:ch04-二项母函数}
\end{align}
\end{example}

\begin{example}[超几何分布的母函数]\label{example:ch04-超几何母函数}
\begin{align}
P(s)=\sum_{k=0}^n\frac{\binom{M}{k}\binom{N-M}{n-k}}{\binom{N}{n}}s^k.\label{align:ch04-超几何母函数}
\end{align}
这是超几何级数，是一种特殊函数，处理起来不大方便，在概率论中也很少用，由于超几何分布的名称来自于它，因此我们顺便提及。
\end{example}

\begin{example}[Poisson 分布的母函数]\label{example:ch04-Poisson母函数}
\begin{align}
P(s)=\sum_{k=0}^{\infty}\frac{\lambda^k}{k!}\mathrm{e}^{-\lambda}s^k
=\mathrm{e}^{-\lambda}\cdot\mathrm{e}^{\lambda s}=\mathrm{e}^{\lambda(s-1)}.\label{align:ch04-Poisson母函数}
\end{align}
\end{example}

\begin{example}[几何分布的母函数]\label{example:ch04-几何母函数}
\begin{align}
P(s)=\sum_{k=1}^{\infty}q^{k-1}ps^k
=ps\sum_{k=1}^{\infty}(qs)^{k-1}=\frac{ps}{1-qs}.\label{align:ch04-几何母函数}
\end{align}
\end{example}

\subsection*{二、母函数的性质}

\textbf{1. 唯一性。}\quad 由分布列~\eqref{align:ch04-母函数定义} 用~\eqref{align:ch04-母函数定义} 定义母函数，这显然是唯一确定的；下面证明，由母函数也能唯一确定分布列。

设概率分布 $\{p_k\}$ 及 $\{q_k\}$ 分别具有母函数 $P(s)$ 及 $Q(s)$，而且 $P(s)=Q(s)$，因为 $P(s)$ 及 $Q(s)$ 都是幂级数，且当 $|s|\leqslant 1$ 时收敛，对 $P(s)$ 及 $Q(s)$ 求导 $k$ 次，并令 $s=0$，则得
\[
k!\,p_k=P^{(k)}(0)=Q^{(k)}(0)=k!\,q_k.
\]
因此 $p_k=q_k$，$k=0,1,2,\cdots$，即两个概率分布一样。

这样一来，概率分布与母函数是一一对应的，因而对于概率分布的许多研究可以化为对所对应的母函数的研究，因为母函数是幂级数，具有许多良好的性质，便于处理，所以母函数是研究整值随机变量的有效工具。

\textbf{2. 母函数与数字特征。}\quad 母函数的应用之一是利用它能求得概率分布的数字特征，若
\[
P(s)=\sum_{k=0}^{\infty}p_k s^k,
\]
即
\[
P'(s)=\sum_{k=1}^{\infty}kp_k s^{k-1},\qquad P''(s)=\sum_{k=2}^{\infty}k(k-1)p_k s^{k-2}.
\]
这两个级数至少在 $|s|<1$ 是收敛的。

当数学期望 $\sum_{k=1}^{\infty}kp_k$ 存在时，
\begin{align}
P'(1)=\sum_{k=1}^{\infty}kp_k=E\xi.\label{align:ch04-母函数求期望}
\end{align}
当数学期望 $\sum_{k=1}^{\infty}kp_k=\infty$ 时，$\lim_{s\to 1}P'(s)=\infty$。

同样，当方差 $D\xi$ 存在时，
\[
E[\xi(\xi-1)]=\sum_{k=2}^{\infty}k(k-1)p_k=P''(1),
\]
故
\begin{align}
D\xi=E\xi^2-(E\xi)^2=P''(1)+P'(1)-[P'(1)]^2.\label{align:ch04-母函数求方差}
\end{align}
公式~\eqref{align:ch04-母函数求期望} 及~\eqref{align:ch04-母函数求方差} 是计算数学期望及方差的简便公式。

\begin{example}[二项分布的数字特征]\label{example:ch04-二项数字特征}
母函数为 $P(s)=(q+ps)^n$。
\[
E\xi=P'(1)=n(q+ps)^{n-1}p\big|_{s=1}=np,
\]
\[
P''(1)=n(n-1)(q+ps)^{n-2}p^2\big|_{s=1}=n(n-1)p^2,
\]
\[
D\xi=n^2p^2-np^2+np-n^2p^2=npq.
\]
这些结果在 $\S 1$ 与 $\S 2$ 中曾直接计算过。
\end{example}

\begin{example}[Poisson 分布的数字特征]\label{example:ch04-Poisson数字特征}
母函数为 $P(s)=\mathrm{e}^{\lambda(s-1)}$。
\[
E\xi=P'(1)=\mathrm{e}^{\lambda(s-1)}\cdot\lambda\big|_{s=1}=\lambda,
\]
\[
P''(1)=\mathrm{e}^{\lambda(s-1)}\cdot\lambda^2\big|_{s=1}=\lambda^2,
\]
\[
D\xi=\lambda^2+\lambda-\lambda^2=\lambda.
\]
这些结果在 $\S 1$ 与 $\S 2$ 中也直接计算过，这里的计算较方便。
\end{example}

\subsection*{三、独立随机变量和的母函数}

若随机变量 $\xi$ 与 $\eta$ 相互独立，它们都是整值随机变量，概率分布分别为 $\{a_k\}$ 及 $\{b_k\}$，而相应的母函数为 $A(s)$ 及 $B(s)$，下面计算随机变量 $\zeta=\xi+\eta$ 的概率分布。显然 $\zeta$ 也是整值随机变量，若记 $c_r=P\{\zeta=r\}$，则
\[
c_r=a_0 b_r+a_1 b_{r-1}+\cdots+a_r b_0.
\]
这就是离散卷积公式~(3.3.6)。

记
\[
C(s)=\sum_{r=0}^{\infty}c_r s^r.
\]
利用母函数在 $|s|\leqslant 1$ 的一致收敛性及绝对收敛性，
\begin{align*}
A(s)B(s)&=\sum_{k=0}^{\infty}a_k s^k\cdot\sum_{l=0}^{\infty}b_l s^l=\sum_{k,l}a_k b_l s^{k+l}\\
&=\sum_{r=0}^{\infty}\left(\sum_{k=0}^r a_k b_{r-k}\right)s^r=\sum_{r=0}^{\infty}c_r s^r.
\end{align*}
因此
\[
C(s)=A(s)B(s).
\]
即两个独立随机变量之和的母函数是这两个随机变量的母函数的乘积。这是一个相当重要的性质，由于母函数具有这个性质，因此在研究独立随机变量和的问题时，母函数很适用。

容易把上面结果推广到 $n$ 个独立整值随机变量之和的场合，若随机变量 $\xi_1,\xi_2,\cdots,\xi_n$ 相互独立，且它们的母函数分别为 $P_1(s),P_2(s),\cdots,P_n(s)$，则 $\eta=\xi_1+\xi_2+\cdots+\xi_n$ 的母函数为
\begin{align}
P(s)=P_1(s)P_2(s)\cdots P_n(s).\label{align:ch04-独立和母函数}
\end{align}

特别当 $\xi_i$ 有相同概率分布的场合，$P_i(s)=P_1(s)$，这时
\begin{align}
P(s)=[P_1(s)]^n.\label{align:ch04-独立同分布和母函数}
\end{align}

\begin{example}[二项分布的母函数再推导]\label{example:ch04-二项母函数再推导}
在成功概率为 $p$ 的 $n$ 次 Bernoulli 试验中，若令
\[
\xi_i=\begin{cases}
1,&\text{在第 }i\text{ 次试验中出现成功},\\
0,&\text{在第 }i\text{ 次试验中出现失败},
\end{cases}
\]
则 $\xi_1,\xi_2,\cdots,\xi_n$ 相互独立，服从 Bernoulli 分布，而 $\eta=\xi_1+\xi_2+\cdots+\xi_n$ 服从二项分布。$\xi_i$ 的母函数为 $(q+ps)$，由~\eqref{align:ch04-独立同分布和母函数}，$\eta$ 的母函数为
\[
P(s)=(q+ps)^n.
\]
这与例 1 直接计算的结果相同。

利用母函数，可以给第二章讲过的 Poisson 逼近定理一个新的证明。

简言之，Poisson 定理可表为：若 $np_n\to\lambda$，则 $B(n,p_n)\to P(\lambda)$。

事实上，$B(n,p_n)$ 的母函数为 $(1-p_n+p_n s)^n$，注意到
\begin{align*}
(1-p_n+p_n s)^n&=\bigl(1+p_n(s-1)\bigr)^n\\
&=\left(1+\frac{np_n(s-1)}{n}\right)^n\to\mathrm{e}^{\lambda(s-1)}.
\end{align*}
这正是 $P(\lambda)$ 的母函数，因此由唯一性推知：$B(n,p_n)\to P(\lambda)$。
\end{example}

\begin{example}[Pascal 分布的母函数]\label{example:ch04-Pascal母函数}
由~(3.1.19) 定义的 Pascal 分布，表示 Bernoulli 试验中第 $r$ 次成功出现时的试验次数 $\zeta$ 的概率分布，而 $\zeta$ 有如下表达式
\[
\zeta=\eta_1+\cdots+\eta_r,
\]
即~(3.1.20)，其中 $\eta_1,\cdots,\eta_r$ 相互独立，均服从几何分布。由~\eqref{align:ch04-几何母函数} 知几何分布的母函数为 $\dfrac{ps}{1-qs}$，因此据~\eqref{align:ch04-独立同分布和母函数} 可得 Pascal 分布的母函数为
\begin{align}
P(s)=\left(\frac{ps}{1-qs}\right)^r.\label{align:ch04-Pascal母函数}
\end{align}

$\eta_1,\cdots,\eta_n$ 的上述两个性质的严格证明，相当繁难，留给有兴趣的读者。
\end{example}

\begin{example}[掷 5 颗骰子]\label{example:ch04-掷骰子}
掷 5 颗骰子，求所得总和为 15 的概率。
\end{example}

\begin{proof}[解]
若以 $\xi_i$ 记第 $i$ 颗掷出的数字，则总和 $\eta=\xi_1+\xi_2+\cdots+\xi_5$，$\xi_i$ 的母函数为
\[
P_i(s)=\frac16(s+s^2+s^3+s^4+s^5+s^6).
\]
显然 $\xi_1,\xi_2,\cdots,\xi_5$ 是相互独立的，因此 $\eta$ 的母函数为
\[
P(s)=\frac{1}{6^5}(s+s^2+\cdots+s^6)^5.
\]
所求的概率 $P\{\eta=15\}$ 是 $P(s)$ 展开式中 $s^{15}$ 项的系数。由于
\begin{align*}
P(s)&=\frac{s^5}{6^5}(1+s+\cdots+s^5)^5=\frac{s^5}{6^5}\left(\frac{1-s^6}{1-s}\right)^5\\
&=\frac{s^5}{6^5}(1-s^6)^5(1-s)^{-5}\\
&=\frac{s^5}{6^5}(1-5s^6+10s^{12}+\cdots-s^{30})\\
&\quad\cdot\left[\sum_{k=0}^{\infty}\binom{-5}{k}(-s)^k\right]\\
&=\frac{s^5}{6^5}(1-5s^6+\cdots)\\
&\quad\cdot\left[\sum_{k=0}^{\infty}(-1)^k\binom{5+k-1}{k}(-1)^k s^k\right]\\
&=\frac{s^5}{6^5}(1-5s^6+\cdots)\left(\sum_{k=0}^{\infty}\binom{k+4}{k}s^k\right).
\end{align*}
故
\[
P\{\eta=15\}=\frac{1}{6^5}\left[1\times\binom{14}{10}-5\times\binom{8}{4}\right]=\frac{651}{6^5}.
\]

掷骰子问题在概率论发展的早期一直占有显著地位，这里用母函数法给予统一处理。
\end{proof}

\subsection*{四、随机个随机变量之和的母函数}

若 $\xi_1,\xi_2,\cdots,\xi_n,\cdots$ 是一串相互独立具有相同概率分布的整值随机变量，$P\{\xi_i=j\}=f_j$，其母函数为
\[
F(s)=\sum_{j=0}^{\infty}f_j s^j.
\]

随机变量 $\nu$ 是取正整数值的，且 $P\{\nu=n\}=g_n$，其母函数为
\[
G(s)=\sum_{n=1}^{\infty}g_n s^n.
\]

若 $\{\xi_n\}$ 与 $\nu$ 独立，考虑和 $\eta=\xi_1+\xi_2+\cdots+\xi_\nu$，记
\[
P\{\eta=i\}=h_i.
\]
我们来求 $\eta$ 的母函数
\[
H(s)=\sum_{i=0}^{\infty}h_i s^i.
\]

利用全概率公式及 $\{\xi_n\}$ 与 $\nu$ 的独立性
\begin{align*}
h_i=P\{\eta=i\}&=\sum_{n=1}^{\infty}P\{\nu=n\}P\{\eta=i\mid\nu=n\}\\
&=\sum_{n=1}^{\infty}P\{\nu=n\}P\{\xi_1+\xi_2+\cdots+\xi_n=i\mid\nu=n\}\\
&=\sum_{n=1}^{\infty}P\{\nu=n\}P\{\xi_1+\xi_2+\cdots+\xi_n=i\}.
\end{align*}

由于 $\xi_1+\xi_2+\cdots+\xi_n$ 为 $n$ 个相互独立具有相同母函数 $F(s)$ 的随机变量之和，故其母函数
\[
\sum_{i=0}^{\infty}P\{\xi_1+\cdots+\xi_n=i\}s^i=[F(s)]^n.
\]
因此
\begin{align}
H(s)&=\sum_{i=0}^{\infty}h_i s^i=\sum_{n=1}^{\infty}P\{\nu=n\}\sum_{i=0}^{\infty}P\{\xi_1+\cdots+\xi_n=i\}s^i\nonumber\\
&=\sum_{n=1}^{\infty}g_n[F(s)]^n=G[F(s)].\label{align:ch04-随机个和母函数}
\end{align}
从~\eqref{align:ch04-随机个和母函数} 可以看到，随机个相互独立相同分布的随机变量之和的母函数是原来两个母函数的复合。

由于
\[
H'(s)=G'[F(s)]\cdot F'(s),
\]
因此当 $E\xi_i$ 及 $E\nu$ 存在时，在上式中令 $s=1$，得到
\begin{align}
E\eta=E\nu\cdot E\xi_i.\label{align:ch04-随机个和期望}
\end{align}
这个公式很直观，也有不少用处。

\textbf{[复合 Poisson 分布]}\quad 在上述讨论中，若 $\nu$ 服从参数为 $\lambda$ 的 Poisson 分布，则
\[
G(s)=\mathrm{e}^{\lambda(s-1)}.
\]
因此 $\eta=\xi_1+\xi_2+\cdots+\xi_\nu$ 的母函数为
\begin{align}
H(s)=\mathrm{e}^{\lambda(F(s)-1)}.\label{align:ch04-复合Poisson母函数}
\end{align}
以~\eqref{align:ch04-复合Poisson母函数} 为母函数的概率分布称为\textbf{复合 Poisson 分布}。这种分布也很有用。

特别当 $F(s)=q+ps$ 时，
\[
H(s)=\mathrm{e}^{\lambda p(s-1)}.
\]

\begin{example}[蚕产卵]\label{example:ch04-蚕产卵}
若每条蚕的产卵数服从 Poisson 分布，参数为 $\lambda$，而每个卵变为成虫的概率为 $p$，且各卵是否变为成虫彼此间没有关系，求每条蚕养活 $k$ 只小蚕的概率。
\end{example}

\begin{proof}[解]
这是第二章习题。现在利用上面结果立刻知道小蚕数服从参数为 $\lambda p$ 的 Poisson 分布，因此所求概率等于
\[
\frac{(\lambda p)^k}{k!}\mathrm{e}^{-\lambda p}.
\]
\end{proof}

\begin{example}[星体数]\label{example:ch04-星体数}
观察资料表明，天空中星体数服从 Poisson 分布，其参数为 $\lambda V$，这里 $V$ 是被观察区域的体积。若每个星球上有生命存在的概率为 $p$，则在体积为 $V$ 的宇宙空间中有生命存在的星球数服从参数为 $\lambda p V$ 的 Poisson 分布。
\end{example}

\textbf{[保险中的索赔模型]}\quad 在非寿命保险中通常采用如下索赔模型：
\begin{align}
\eta=\xi_1+\xi_2+\cdots+\xi_\mu,\label{align:ch04-索赔模型}
\end{align}
这里 $\xi_i$，$i=1,2,\cdots$ 为独立同分布随机变量，表示第 $i$ 次索赔数额，而总索赔次数则是随机变量 $\mu$，因此表现为随机个独立随机变量之和的形式。更一般的索赔模型则还假定 $\mu$ 也是随机个随机变量之和：
\begin{align}
\mu=\theta_1+\theta_2+\cdots+\theta_\nu.\label{align:ch04-索赔次数模型}
\end{align}

\end{document}